\section{Conclusion}
\label{sec:conclusion}

We have proposed PY-IDR, a hierarchical nonparametric Bayesian extension of the IDR framework for multi-replicate reproducibility analysis. The model combines a Pitman--Yor mixture for reproducible-component dependence regimes, structured Gaussian-copula correlation classes augmented with planned heavy-tail atoms, and replicate-specific Bernstein--Dirichlet marginal priors. We developed a P\'olya-urn auxiliary-atom MCMC scheme and a finite-truncation variational approximation, and we stated theoretical results with explicit assumptions separating finite-sieve identifiability, conditional contraction, threshold-regularity-based Bayes-mFDR control, Harris ergodicity, and conjectural or conditional stronger claims. The manuscript also lays out a protocol-first empirical program with simulations, ablations, real-data analyses, placeholder figures, and placeholder tables to be completed in the companion experimental implementation.

\subsection{Limitations}

Several limitations are central to the current manuscript status. First, the identifiability result is a finite-sieve statement with explicit linear-independence and component-separation assumptions; extending it to unrestricted infinite PY mixtures requires additional weak-identifiability theory for the mixing measure. Second, the contraction theorem is conditional on a sieve-tail assumption. For positive Pitman--Yor discount, the polynomial residual expectation is not by itself an exponential sieve-complement bound, so this remains a proof obligation for a fully unconditional PY-copula contraction theorem. Third, the stated contraction theory covers bounded smooth copula truths; boundary-singular tail-dependent copulas motivate the methodology and S5 experiments but require additional boundary theory. Fourth, Bayes-mFDR rate statements require local-idr error bounds and threshold regularity beyond global Hellinger contraction alone. Fifth, geometric ergodicity and variational matching-rate claims remain conditional on additional drift, small-set, truncation, and variational-approximation assumptions.

\subsection{Practical recommendations}

For practitioners using the current PY-IDR design, recommendations should be treated as provisional implementation guidance until the companion experiments are run. MCMC is the preferred default for moderate $n$ when posterior uncertainty and calibration are primary, whereas VI is intended for large genomic-scale runs where approximate posterior summaries are acceptable. The auxiliary-atom count $H=5$ is a reasonable starting value for the planned sensitivity grid, not an empirically validated optimum. Heavy-tail atoms should be included in the planned ChIP-seq, ATAC-seq, and GWAS stress tests, but their final effect on discoveries must be assessed by the completed ablation study. Dataset-specific rank handling, tie handling, missingness policy, and feature-universe construction should be frozen before model fitting.

\subsection{Future work}

Several extensions are natural. The most direct extension is to vine-copula atoms in place of structured-correlation Gaussian atoms, which would allow asymmetric pairwise dependence patterns at the cost of a more demanding base measure. A second extension is to time-varying or covariate-dependent dependence regimes, in the framework of \citet{DallaValleLeisenRossiniZhu2024}, which would accommodate longitudinal reproducibility studies. A third extension is to non-exchangeable replicate structures via dependent Pitman--Yor priors, which would allow modeling of donor-stratified single-cell data with within-donor and between-donor dependence regimes. A fourth extension is to multi-resolution score representations beyond ranks, exploiting original signal magnitudes through $z$-score-aware likelihoods. Finally, a complete drift-and-minorization proof for the P\'olya-urn PY-IDR sampler and an unconditional contraction theorem for positive-discount PY copula mixtures would substantially strengthen the theoretical contribution.

\subsection*{Acknowledgments and code availability}

\paragraph{Acknowledgments.}
[Acknowledgments to be added in the final version.]

\paragraph{Funding.}
[Funding sources to be added in the final version.]

\paragraph{Code availability.}
A reference implementation is planned as an open-source R/Python package, with reproducible workflow scripts for the real-data applications. The GitHub repository, archival DOI, and frozen container image will be linked only after the implementation and companion experiments are complete.

\paragraph{Data availability.}
The planned datasets are publicly available through the cited consortia or resources, subject to each resource's access terms and data-use requirements. Dataset manifests and accession identifiers will be reported in the companion experimental paper.

\paragraph{Conflict of interest.}
The authors declare no competing interests.
